\documentclass[11pt,a4paper]{article}
\usepackage[margin=25mm]{geometry}
\usepackage{amsmath,amssymb,amsthm}
\newtheorem{lemma}{Lemma}\newtheorem{proposition}[lemma]{Proposition}
\newcommand{\Sz}{\operatorname{Sz}}
\title{Cyclic overlaps of two Suzuki subfield faces}
\author{Human Compression project: bounded F003 continuation}\date{}
\begin{document}\maketitle
This is a bounded written reduction of the remaining case in the accepted
\emph{The remaining subfield obstruction for Suzuki rotation markings}.
The entire two-page proper-face paper and entire two-page subfield paper
remain written premises; the exact bindings and classical inputs are
listed in the companion dependency ledger. This reduction does not close
that case. In the predecessor notation let
$F=\langle a,b\rangle\cong\Sz(2^d)$ and
$E=\langle b,c\rangle\cong\Sz(2^e)$ be proper subfield groups of
$S=\Sz(2^N)$, with $F\cap E=\langle b\rangle$ and $b^2\ne1$.
All three exponents are odd and $d,e>1$ divide $N$.
We use the retained classical standard subfield, ovoid and torus
interfaces, not an unproved equality of differently embedded subfields.
Put $h=\gcd(d,e)$.

\begin{lemma}[Root-centre overlap]\label{root}
If $|b|=4$, then the centres of the point Sylow $2$-groups of $F$ and
$E$ containing $b$ intersect in a group of order $2^h$.
\end{lemma}
\begin{proof}
The unique ambient fixed point $P$ of $b$ lies in both subovoids, by
fixed-point descent in the paid predecessor. Conjugate $F$ to the
standard subgroup over $K=\mathbb F_{2^d}$ and $P$ to infinity within
its subovoid. Write $L=\mathbb F_{2^e}$, $Q=2^N$ and
$\theta=2^{(N+1)/2}$.
The standard coordinates, from Bray--B\"a\"arnhielm (2017), equations
(1)--(5), are
\[
 U(x,y)=\begin{pmatrix}
 1&0&0&0\\x&1&0&0\\x^{\theta+1}+y&x^\theta&1&0\\
 x^{\theta+2}+xy+y^\theta&y&x&1
 \end{pmatrix},\qquad
 H(\lambda)=\operatorname{diag}(\lambda^{\theta+1},\lambda,
 \lambda^{-1},\lambda^{-\theta-1}).
\]
Matrix multiplication shows that the first coordinate is additive
under root-group multiplication. Conjugation by any root-group element
therefore preserves it; such conjugation also fixes every central
$U(0,y)$. Conjugation by $H(\lambda)$ scales the first coordinate by
$\mu=\lambda^\theta$ and the central coordinate by
$\lambda^{\theta+2}=\mu^{\theta+1}$. These identities hold on $K,L$
as well because the ambient twist restricts to each subfield twist.

A conjugator taking the standard $\Sz(2^e)$ to $E$ may be chosen in the
ambient stabilizer of $P$: adjust any conjugator by an element of the
standard smaller group taking the inverse image of $P$ to its infinity.
Its root-group part leaves the first-coordinate set and centre
unchanged. Thus the first-coordinate sets of the two point root groups
are $K$ and $\mu L$, and their central-coordinate sets are $K$ and
$\mu^{\theta+1}L$, for some $\mu\ne0$.

Write $b=U(x,y)$. Direct multiplication gives
$U(x,y)^2=U(0,x^{\theta+1})$. Order four therefore implies $x\ne0$;
the $U(0,y)$ are precisely the central involutions and identity. The common membership
gives $x\in K\cap\mu L$. Write $x=\mu\ell$ with
$\ell\in L^\times$. Then
\[
 K\cap\mu^{\theta+1}L
 =K\cap x^{\theta+1}L
 =x^{\theta+1}(K\cap L).
\]
The last equality uses $x^{\theta+1}\in K^\times$. Since
$K\cap L=\mathbb F_{2^h}$, the claimed order follows.
\end{proof}

\begin{proposition}[Restrictions on the middle rotation]\label{middle}
If $|b|=4$, then $\gcd(d,e)=1$. If $b$ has odd order, it generates
the intersection of its two face centralizers:
\[
 \langle b\rangle=C_F(b)\cap C_E(b),\qquad
 |b|=\gcd(|C_F(b)|,|C_E(b)|).
\]
In particular, if $\gcd(d,e)=1$, then $|b|\in\{4,5\}$.
\end{proposition}
\begin{proof}
In the order-four case the elementary abelian group of
Lemma~\ref{root} is contained in $F\cap E=\langle b\rangle$. A cyclic
group of order four contains just two elements of order dividing two.
Therefore $2^h\le2$, so $h=1$.

For odd $b$, its unique ambient centralizer is cyclic. Both face
centralizers are subgroups of this one cyclic group, and their
intersection contains $\langle b\rangle$ and is contained in
$F\cap E$. This proves the equality and the gcd formula.

Assume now $h=1$. Split and nonsplit types cannot be mixed: a common
odd divisor of $2^d-1$ and $2^{2e}+1$ is one. Indeed modulo such a
divisor the order of two is both odd (dividing $d$) and divides $4e$,
so it divides $\gcd(d,e)=1$; the second congruence then makes the
divisor divide two. The same proof works with $d,e$ interchanged.
If both types are split, the gcd is $2^h-1=1$, also impossible.
If both are nonsplit, their orders divide $2^{2d}+1$ and $2^{2e}+1$.
For coprime odd $d,e$,
\[
 \gcd(2^{2d}+1,2^{2e}+1)=2^2+1=5.
\]
For completeness, if $D$ is the gcd and $A=4$, Bezout gives integers
$r,s$ with $rd+se=1$. Their sum is odd, because $d,e$ are odd.
The congruences $A^d=A^e=-1\pmod D$ give $A=-1\pmod D$, hence
$D\mid5$; conversely both displayed numbers are divisible by five.
Thus a nonidentity odd $b$ has order five. The retained Suzuki element
types say that every nonidentity element has odd order or order two or
four; $b^2\ne1$ excludes the remaining even possibilities.
\end{proof}

In particular, equal proper subfield degrees exclude an order-four
middle rotation. This does not exclude an odd middle rotation, or the
coprime-degree order-four/order-five alternatives. The trace probes
for $\Sz(8)$ and $\Sz(32)$ fit the latter boundary, but their absence
results are not premises here and do not establish a uniform theorem.
\end{document}
